\section{Auxiliary estimates}\label{sec:auxiliary}

Before entering the two branches of the proof, we collect the three main
tools: bounded probe functions for reading off coefficients from the
quadratic observable $X_a^2 - X_b^2$; an $L^1$ lower bound for independent
sums; and a gluing estimate showing that if $a$ and $\tau b$ nearly agree
\emph{everywhere} (head and tail), then $\eps$ is already large.

\subsection{Probe functions}\label{sec:probes}

Given a centered random variable $\xi$ with $\E[\abs{\xi}] > 0$, one can
always find a bounded function of $\xi$ that acts as a dual to $\xi$ itself.
If additionally $\xi$ is not Rademacher-valued (i.e., $\E[\abs{\xi}] < 1$),
the same construction works for $\xi^2 - 1$. These are the probes that let
us read off individual entries of $D_H$ from $\eps$.

\begin{lemma}[Linear probe]\label{lem:linear-probe}
\leanname{Imported\.ElementaryBounds\.exists\_linear\_probe}
Let $\xi$ be a real random variable with $\E[\xi] = 0$ and
$\E[\abs{\xi}] \ge A > 0$. Then there exists a measurable
$g : \R \to \R$ with $\abs{g} \le 2/A$,\, $\E[g(\xi)] = 0$, and
$\E[\xi \, g(\xi)] = 1$.
\end{lemma}

\begin{proof}
Let $D = \E[\abs{\xi}]$ and $c = \E[\sign(\xi)]$, and set
$g(x) = (\sign(x) - c)/D$.
Then $\E[g(\xi)] = 0$ and
$\E[\xi \, g(\xi)] = \E[\abs{\xi}]/D = 1$. The bound $\abs{g} \le 2/D$
follows from $\abs{\sign(x) - c} \le 2$.
\end{proof}

\begin{lemma}[Quadratic probe]\label{lem:quadratic-probe}
\leanname{Imported\.ElementaryBounds\.exists\_quadratic\_probe}
Let $\xi$ be a real random variable with $\E[\xi] = 0$,
$\E[\xi^2] = 1$, and $\E[\abs{\xi}] \le B < 1$. Then there exists a
measurable $g : \R \to \R$ with $\abs{g} \le 2/(1-B)$,\,
$\E[g(\xi)] = 0$, and $\E[(\xi^2 - 1) \, g(\xi)] = 1$.
\end{lemma}

\begin{proof}
Apply the linear-probe construction to $\eta = \xi^2 - 1$, whose $L^1$ norm
satisfies
$\E[\abs{\xi^2 - 1}] \ge 1 - \E[\abs{\xi}] \ge 1 - B > 0$.
\end{proof}

\begin{remark}\label{rem:probe-extraction}
The linear probe recovers $c_i$ from $X_c = \sum_k c_k \xi_k$ via the
identity $\E[X_c \cdot g(\xi_i)] = c_i$: the centering $\E[g(\xi_i)] = 0$
kills the contribution of every $\xi_k$ with $k \ne i$, and the
normalization $\E[\xi_i g(\xi_i)] = 1$ picks out $c_i$. Likewise, the
quadratic probe gives $\E[X_c^2 \cdot g(\xi_i)] = c_i^2$.
\end{remark}

\subsection{The $L^1$ lower bound}\label{sec:l1-lower}

Independent sums of variables with positive $L^1$ norms cannot concentrate
near zero. This is the engine behind both the gluing estimate and the
diffuse branch: it converts $\ell^2$ energy of the coefficients into $L^1$
mass of the random sum.

\begin{theorem}[$L^1$ lower bound for independent sums]%
\label{thm:l1-lower}
\leanname{Imported\.ElementaryBounds\.finite\_indep\_sum\_l1\_lower}
Let $\zeta_1, \dots, \zeta_n$ be independent, mean-zero, unit-variance
random variables with $\E[\abs{\zeta_i}] \ge A > 0$ for all $i$. Then
for every $(c_1, \dots, c_n) \in \R^n$,
\[
\E\Bigl[\Bigl\lvert \sum_{i=1}^n c_i \zeta_i
\Bigr\rvert\Bigr]
\ge \frac{A}{4}
\sqrt{\sum_{i=1}^n c_i^2}\,.
\]
\end{theorem}

\subsection{The four-term gluing estimate}\label{sec:gluing}

Orthogonality forces $\norm{a + \tau b}_2 = \norm{a - \tau b}_2 = \sqrt{2}$.
If for some $\tau$ both the head of $a - \tau b$ and the tail of
$a + \tau b$ are small, then the complementary pieces---the tail of
$a - \tau b$ and the head of $a + \tau b$---must be close to $\sqrt{2}$.
The product $X_{a+\tau b} \cdot X_{a-\tau b} = X_a^2 - X_b^2$ then has a
large main term and small errors, yielding a lower bound on $\eps$.

\begin{lemma}[Four-term bound]\label{lem:four-term}
\privateleanname{Tail\.SignMismatch\.fin\_sign\_mismatch\_measure\_bound}
Let $a, b \in \R^n$ be unit vectors with $\ip{a}{b} = 0$, let
$H \subseteq \{1, \dots, n\}$, and let $\tau \in \{+1, -1\}$. Write
$c^+ = a + \tau b$ and $c^- = a - \tau b$. Then
\begin{align*}
\eps(a,b) \;\ge\;
&\bigl(\tfrac{A}{4}\bigr)^2\,
  \norm{\restr_H c^+}_2 \, \norm{\tail_H c^-}_2 \\
&{}- \norm{\restr_H c^+}_2 \, \norm{\restr_H c^-}_2
 - \norm{\tail_H c^+}_2 \, \norm{\tail_H c^-}_2 \\
&{}- \norm{\tail_H c^+}_2 \, \norm{\restr_H c^-}_2.
\end{align*}
\end{lemma}

\begin{proof}
The identity $X_a^2 - X_b^2 = X_{c^+} \cdot X_{c^-}$ holds since
$\tau^2 = 1$. Set $P = X_{\restr_H c^+}$, $X' = X_{\tail_H c^+}$,
$W = X_{\restr_H c^-}$, $V = X_{\tail_H c^-}$. Expanding
$(P + X')(W + V)$ and rearranging:
\[
\abs{PV}
\le \abs{X_{c^+} X_{c^-}} + \abs{PW} + \abs{X'V} + \abs{X'W}.
\]
Taking expectations, the left-hand side factors as
$\E[\abs{P}] \cdot \E[\abs{V}]$ by independence of disjoint supports
(\leanref{Imported\.L2Theory\.indepFun\_finLinComb\_disjoint}).
Each factor is at least $(A/4)$ times the $\ell^2$ norm of the
corresponding coefficient vector, by Theorem~\ref{thm:l1-lower}. The
three error terms are each bounded by the product of the relevant $L^2$
norms via Cauchy--Schwarz.
\end{proof}

\begin{proposition}[Gluing]\label{thm:gluing}
\leanname{Tail\.SignMismatch\.fin\_sign\_mismatch\_gluing}
Let $a, b \in \R^n$ be unit vectors with $\ip{a}{b} = 0$, let
$H \subseteq \{1, \dots, n\}$, and let $\tau \in \{+1, -1\}$.
If
\[
\norm{\restr_H(a - \tau b)}_2 \le \theta(A)/16
\quad \text{and} \quad
\norm{\tail_H(a + \tau b)}_2 \le \theta(A)/32,
\]
then $\eps(a,b) \ge \theta(A)/2$.
\end{proposition}

\begin{proof}
Write $c^+ = a + \tau b$ and $c^- = a - \tau b$. Set
$\alpha = \norm{\restr_H c^-}_2$ and $\beta = \norm{\tail_H c^+}_2$;
by hypothesis $\alpha \le \theta/16$ and $\beta \le \theta/32$.
From $\norm{c^\pm}_2^2 = 2$ we get
\[
\norm{\restr_H c^+}_2
= \sqrt{2 - \beta^2} \ge 1,
\qquad
\norm{\tail_H c^-}_2
= \sqrt{2 - \alpha^2} \ge 1.
\]
The main term in the four-term bound (Lemma~\ref{lem:four-term}) is
therefore at least
$(A/4)^2 \cdot 1 \cdot 1 = \theta/2$, while the three error terms sum to
at most $\alpha \sqrt{2} + \beta \sqrt{2} + \alpha\beta < \theta/4$.
The conclusion $\eps \ge \theta/2$ follows.
\end{proof}
